% problem-000330 2 液氨溶剂与金属溶液
\section*{2. 液氨溶剂与金属溶液}
\textit{下列为分问。}

\subsection*{2-1}
作为溶剂， $\mathrm { N H _ { 3 } }$ 分⼦也能发⽣类似于 $\mathrm { H _ { 2 } O }$ 分⼦的缔合作⽤，说明发⽣这种缔合的原因和⽐较这种缔合作⽤相对于⽔的⼤⼩。

\subsection*{2-2}
以液氨作为溶剂最引起化学家兴趣的是它能够溶解⼀些⾦属，如电极电势⼩于2.5V的碱⾦属、部分碱⼟⾦属及镧系元素可溶于液氨，形成蓝⾊的具有异乎寻常性质的亚稳定态溶液，这种溶液具有顺磁性和⾼的导电性，溶液的密度⽐纯溶液的密度⼩。碱⾦属的液氨溶液是可供选择使⽤的强还原剂，⼴泛应⽤于合成⼀些⾮正常氧化态的⾦属配合物和其他化合物。研究发现在⾦属Na的液氨溶液中存在着以下反应：

$
\begin{array}{l}\mathrm {Na(s) + NH_ {3} \rightarrow Na(NH_ {3})}\\2 \mathrm {Na(NH_ {3}) \rightleftharpoons Na_ {2} (NH_ {3}) + NH_ {3}}\end{array}\quad K ^ {\ominus} \approx 5 \times 1 0 ^ {3}
$

$
\begin{array}{l l}\mathrm {Na(NH_ {3}) + NH_ {3} \rightleftharpoons Na^ {+} (NH_ {3}) + e^ {-} (NH_ {3})}&K ^ {\ominus} \approx 1 0 ^ {- 2}\\\mathrm {Na^ {-} (NH_ {3}) + NH_ {3} \rightleftharpoons Na(NH_ {3}) + e^ {-} (NH_ {3})}&K ^ {\ominus} \approx 1 0 ^ {- 3}\\2 \mathrm {e^ {-} (NH_ {3}) \rightleftharpoons e_ {2} ^ {2 - } (NH_ {3}) + NH_ {3}}\end{array}
$

⾦属 $\mathrm { N a }$ 在液 $\mathrm { N H _ { 3 } }$ 溶剂中⽣成氨合⾦属 $\mathrm { N a } ( \mathrm { N H } _ { 3 } )$ $\mathrm { N a _ { 2 } ( N H _ { 3 } ) }$ ，氨合阳离⼦ $\cdot \mathrm { N a ^ { + } ( N H _ { 3 } ) }$ ，氨合阴离⼦ $\mathrm { N a ^ { - } ( N H _ { 3 } ) }$ 及氨合电⼦e−(NH_{3})、 $\mathrm { e } _ { 2 } ^ { 2 \mathrm { - } } \mathrm { ( N H _ { 3 } ) }$ 等物种。

试根据以上信息解释碱⾦属液氨溶液的⾼导电性和顺磁性。

出 $\mathrm { | [ P t ( N H _ { 3 } ) _ { 4 } ] B r _ { 2 } }$ 与K在液氨中的反应⽅程式。

\subsection*{2-3}
溶剂 $\mathrm { N H _ { 3 } }$ 分⼦⾃动电离形成铵离⼦和氨基离⼦：

$
2 \mathrm{NH} _ {3} \rightleftharpoons \mathrm{NH} _ {4} ^ {+} + \mathrm{NH} _ {2} ^ {-} \quad K ^ {\ominus} = [ \mathrm{NH} _ {4} ^ {+} ] \times [ \mathrm{NH} _ {2} ^ {-} ] = 1 0 ^ {- 2 7}
$

离⼦积常数 $K ^ { \ominus }$ 虽然⽐ $\mathrm { H _ { 2 } O }$ 的 $K _ { \mathbf { w } } ^ { \ominus }$ ⼩得多，但同样可以建⽴类似于⽔体系中的 $\mathrm { p H }$ 标度。试建⽴这种标度，确定酸性、中性和碱性溶液的 $\mathrm { p H }$ C

\subsection*{2-4}
在液氨中的许多反应都类似于⽔中的反应，试写出 $\mathrm { T i C l _ { 4 } , \ Z n ^ { 2 + } }$ 和 $\mathrm { L i _ { 3 } N }$ 在液氨中反应的⽅程式。

\subsection*{2-5}
在⽔体系中以标准氢电极为基准建⽴了标准电极电势系统，在液氨体系中同样也以标准氢电极建⽴类似于⽔体系的液氨体系中的标准电极电势系统。试写出标准氢电极的半反应⽅程式及标准电极电势。

\subsection*{2-6}
在液氨体系中某些⾦属的标准电极电势跟这些⾦属在⽔溶液的标准电极电势⼗分相近。已知在液氨体系中， $6 \mathrm { N H _ { 4 } ^ { + } + N _ { 2 } + 6 e ^ { - }  8 N H _ { 3 } }$ 的 $E ^ { \ominus } = 0 . 0 4 \mathrm { V } _ { \odot }$ 。试设计⼀个在液氨中实施⽤氢固氮的反应，并预计其反应条件，简述理由。

\subsection*{2-7}
液氨作为溶剂在化学分析中⼴泛⽤于⾮⽔滴定，试述哪些不宜于在⽔溶液中滴定的酸碱体系可以在液氨中进⾏。
